\documentclass[preprint,12pt]{elsarticle}

\usepackage{amsmath,amssymb,amsfonts,amsthm}
\usepackage{array}
\usepackage{booktabs}
\usepackage{graphicx}
\usepackage{url}
\usepackage{xcolor}
\providecommand{\nolinkurl}[1]{\url{#1}}

\journal{Computer Methods in Applied Mechanics and Engineering}
\emergencystretch=4em

\newtheorem{assumption}{Assumption}
\newtheorem{definition}{Definition}
\newtheorem{lemma}{Lemma}
\newtheorem{theorem}{Theorem}
\newtheorem{proposition}{Proposition}
\newtheorem{corollary}{Corollary}

\newcommand{\method}{Gauss6\slash FullVA}
\newcommand{\tfe}{TFE}
\newcommand{\figpath}{figures}
\newcolumntype{P}[1]{>{\raggedright\arraybackslash}p{#1}}
\renewcommand{\topfraction}{0.9}
\renewcommand{\bottomfraction}{0.7}
\renewcommand{\textfraction}{0.08}
\renewcommand{\floatpagefraction}{0.85}
\setcounter{topnumber}{3}
\setcounter{totalnumber}{4}
\newcommand{\paperfig}[3]{%
  \begin{figure}[!htb]
    \centering
    \IfFileExists{#1}{%
      \includegraphics[width=\linewidth]{#1}%
    }{%
      \fbox{\parbox{0.86\linewidth}{Missing figure: \nolinkurl{#1}}}%
    }
    \caption{#2}
    \label{#3}
  \end{figure}
}

\begin{document}

\begin{frontmatter}

\title{A sixth-order Lie-group Gauss collocation integrator for lower-pair mechanisms with position-, velocity- and acceleration-level constraints}

\author[inst1]{Jingquan Wang\corref{cor1}}
\ead{jingquanw@wisc.edu}
\cortext[cor1]{Corresponding author.}

\affiliation[inst1]{organization={University of Wisconsin--Madison},
  addressline={Department of Mechanical Engineering, 1513 University Avenue},
  city={Madison},
  state={WI},
  postcode={53706},
  country={United States}}

\begin{abstract}
Multibody systems with lower-pair joints in absolute coordinates are index-3 differential-algebraic
equations on a Lie group. We study \method{}, a three-stage Gauss collocation step whose stage
system enforces the joint constraints at position, velocity and acceleration level together with
the Newton--Euler balance, and reconstructs the rotational endpoint through the exponential map.
The central result is an exact algebraic identity: the lifted reduced Gauss stage, obtained by
collocating the joint-coordinate equations of motion and lifting the result to absolute
coordinates, satisfies every one of the $132$ implemented stage rows, and conversely every root of
the non-dynamic rows on the regular branch is such a lift. The method is therefore reduced Gauss
collocation executed in absolute coordinates. Everything the implementation adds, the
velocity-level endpoint closure and the inexact Newton solve, is an $O(h^7)$ perturbation, which
gives a sixth-order error bound whose only standing hypothesis is smoothness on a compact
neighbourhood of the trajectory; the solver and inverse interfaces are derived for $h\le h_0$, and
$h_0$ is quantified and traced to the curvature of the friction law. The row identities, the
perturbation chain and the Butcher conditions are machine-checked in Lean~4 with Mathlib
($57$ theorems). Numerically the method is sixth order on a frictional two-body chain and on the
ASME double pendulum, reproduces the driven ASME mechanisms to roundoff, is two to seven orders of
magnitude more accurate than the one- and two-stage members of the same family at the same Newton
work, and reaches errors below $10^{-10}$ on the double pendulum where the public
absolute-coordinate codes have errors of order one. A sweep of the Stribeck velocity locates the
regime in which a sharp friction law reduces the observed order on coarse grids.
\end{abstract}

\begin{keyword}
Lie group integrator \sep index-3 DAE \sep multibody dynamics \sep
lower-pair mechanisms \sep Gauss collocation \sep machine-checked proof
\end{keyword}

\end{frontmatter}
